
\documentclass[12pt]{article}
\usepackage{math_hw}
\usepackage{amsmath, amssymb}

% The true underlying data generating distribution
\newcommand{\pdata}{p_{\rm{data}}}
% The empirical distribution defined by the training set
\newcommand{\ptrain}{\hat{p}_{\rm{data}}}
\newcommand{\Ptrain}{\hat{P}_{\rm{data}}}
% The model distribution
\newcommand{\pmodel}{p_{\rm{model}}}
\newcommand{\Pmodel}{P_{\rm{model}}}
\newcommand{\ptildemodel}{\tilde{p}_{\rm{model}}}
% Stochastic autoencoder distributions
\newcommand{\pencode}{p_{\rm{encoder}}}
\newcommand{\pdecode}{p_{\rm{decoder}}}
\newcommand{\precons}{p_{\rm{reconstruct}}}

\newcommand{\laplace}{\mathrm{Laplace}} % Laplace distribution

\newcommand{\E}{\mathbb{E}}
\newcommand{\Ls}{\mathcal{L}}
\newcommand{\R}{\mathbb{R}}
\newcommand{\emp}{\tilde{p}}
\newcommand{\lr}{\alpha}
\newcommand{\reg}{\lambda}
\newcommand{\rect}{\mathrm{rectifier}}
\newcommand{\softmax}{\mathrm{softmax}}
\newcommand{\sigmoid}{\sigma}
\newcommand{\softplus}{\zeta}
\newcommand{\KL}{D_{\mathrm{KL}}}
\newcommand{\Var}{\mathrm{Var}}
\newcommand{\standarderror}{\mathrm{SE}}
\newcommand{\Cov}{\mathrm{Cov}}
% Wolfram Mathworld says $L^2$ is for function spaces and $\ell^2$ is for vectors
% But then they seem to use $L^2$ for vectors throughout the site, and so does
% wikipedia.
\newcommand{\normlzero}{L^0}
\newcommand{\normlone}{L^1}
\newcommand{\normltwo}{L^2}
\newcommand{\normlp}{L^p}
\newcommand{\normmax}{L^\infty}

\newcommand{\xpos}{x_+}
\newcommand{\xneg}{x_-}

\newcommand{\parents}{Pa} % See usage in notation.tex. Chosen to match Daphne's book.

\DeclareMathOperator*{\argmax}{arg\,max}
\DeclareMathOperator*{\argmin}{arg\,min}

\DeclareMathOperator{\sign}{sign}
\let\ab\allowbreak
\newcommand{\ones}{\mathbf 1}
\newcommand{\reals}{\mathbb{R}}
\newcommand{\integers}{{\mbox{\bf Z}}}
\newcommand{\symm}{{\mbox{\bf S}}}  % symmetric matrices
\DeclareMathOperator*{\esssup}{ess\,sup}
\newcommand{\nullspace}{{\mathcal N}}
\newcommand{\range}{{\mathcal R}}
\newcommand{\diag}{\mathop{\bf diag}}
\newcommand{\card}{\mathop{\bf card}}
\newcommand{\conv}{\mathop{\bf conv}}
\newcommand{\prox}{\mathbf{Prox}}

\newcommand{\Expect}{\mathop{\bf E{}}}
\newcommand{\Prob}{\mathop{\bf Prob}}
\newcommand{\Co}{{\mathop {\bf Co}}} % convex hull
\newcommand{\dist}{\mathop{\bf dist{}}}
\newcommand{\epi}{\mathop{\bf epi}} % epigraph
\newcommand{\Vol}{\mathop{\bf vol}}
\newcommand{\dom}{\mathop{\bf dom}} % domain
\newcommand{\intr}{\mathop{\bf int}}
\newcommand{\avar}{\mathbf{AVar}}
\newcommand{\expec}{\mathbb{E}}

\newcommand{\cf}{{\it cf.}}
\newcommand{\eg}{{\it e.g.}}
\newcommand{\ie}{{\it i.e.}}
\newcommand{\etc}{{\it etc.}}


\newcommand{\fC}{\field{C}} 
\newcommand{\fE}{\field{E}} 
\newcommand{\fN}{\field{N}} % natural numbers
\newcommand{\fR}{\field{R}} % real numbers
\newcommand{\fZ}{\field{Z}} % integers

\newcommand{\bzero}[1]{0}



\newcommand{\mathify}[1]{\ifmmode{#1}\else\mbox{$#1$}\fi}
\newcommand{\Epi}[1]{\mathify{\mathbb{E}_{\pi}\left[#1\right]}}	
\newcommand{\Ep}[1]{\mathify{\mathbb{E}_{P}\left[#1\right]}}
\newcommand{\Eq}[1]{\mathify{\mathbb{E}_{Q}\left[#1\right]}}
\newcommand{\EX}[1]{\mathify{\mathbb{E}_{X}\left[#1\right]}}
\newcommand{\EZ}[1]{\mathify{\mathbb{E}_{Z}\left[#1\right]}}




\newcommand\blfootnote[1]{%
  \begingroup
  \renewcommand\thefootnote{}\footnote{#1}%
  \addtocounter{footnote}{-1}%
  \endgroup
}
%%%%% NEW MATH DEFINITIONS %%%%%

\usepackage{amsmath,amsfonts,bm}

% Mark sections of captions for referring to divisions of figures
\newcommand{\figleft}{{\em (Left)}}
\newcommand{\figcenter}{{\em (Center)}}
\newcommand{\figright}{{\em (Right)}}
\newcommand{\figtop}{{\em (Top)}}
\newcommand{\figbottom}{{\em (Bottom)}}
\newcommand{\captiona}{{\em (a)}}
\newcommand{\captionb}{{\em (b)}}
\newcommand{\captionc}{{\em (c)}}
\newcommand{\captiond}{{\em (d)}}

% Highlight a newly defined term
\newcommand{\newterm}[1]{{\bf #1}}

\newcommand{\field}[1]{\mathbb{#1}}
\newcommand{\mse}{\text{\rm mse}}
\newcommand{\Ascr}{{\cal A}}
\newcommand{\Bscr}{{\cal B}}
\newcommand{\Cscr}{{\cal C}}
\newcommand{\Dscr}{{\cal D}}
\newcommand{\Escr}{{\cal E}}
\newcommand{\Fscr}{{\cal F}}
\newcommand{\Gscr}{{\cal G}}
\newcommand{\Hscr}{{\cal H}}
\newcommand{\Iscr}{{\cal I}}
\newcommand{\Jscr}{{\cal J}}
\newcommand{\Kscr}{{\cal K}}
\newcommand{\Lscr}{{\cal L}}
\newcommand{\Mscr}{{\cal M}}
\newcommand{\Nscr}{{\cal N}}
\newcommand{\Oscr}{{\cal O}}
\newcommand{\Pscr}{{\cal P}}
\newcommand{\Qscr}{{\cal Q}}
\newcommand{\Rscr}{{\cal R}}
\newcommand{\Sscr}{{\cal S}}
\newcommand{\Tscr}{{\cal T}}
\newcommand{\Uscr}{{\cal U}}
\newcommand{\Vscr}{{\cal V}}
\newcommand{\Wscr}{{\cal W}}
\newcommand{\Xscr}{{\cal X}}
\newcommand{\Yscr}{{\cal Y}}
\newcommand{\Zscr}{{\cal Z}}

\newcommand{\Xb}[1]{{\bf X^{(#1)}}}
\newcommand{\Bb}[1]{{\bf B^{(#1)}}}
\newcommand{\Xe}[1]{ X^{(#1)}}
\newcommand{\Be}[1]{ B^{(#1)}}
\newcommand{\Zb}[1]{{\bf Z^{(#1)}}}
\newcommand{\Ze}[1]{ Z^{(#1)}}
\newcommand{\Yb}[1]{{\bf Y^{(#1)}}}
\newcommand{\Ye}[1]{ Y^{(#1)}}
\newcommand{\Ub}[1]{{\bf U^{(#1)}}}
\newcommand{\Ue}[1]{ U^{(#1)}}
\newcommand{\Vb}[1]{{\bf V^{(#1)}}}
\newcommand{\Ve}[1]{ V^{(#1)}}
\newcommand{\Wb}[1]{{\bf W^{(#1)}}}
\newcommand{\We}[1]{ W^{(#1)}}
\newcommand{\Wbtilde}[1]{{\bf \widetilde{W}^{(#1)}}}
\newcommand{\Wetilde}[1]{ \widetilde{W}^{(#1)}}
\newcommand{\Ab}{{\bf A}}
\newcommand{\Qb}{{\bf Q}}
\newcommand{\Xbhat}{{\bf \hat{X}}}
\newcommand{\Xbh}[1]{{\bf \hat{X}^{(#1)}}}
\newcommand{\mb}[1]{{\bf m^{(#1)}}}
\newcommand{\me}[1]{ m^{(#1)}}
\newcommand{\mbhat}[1]{{\bf \hat{m}^{(#1)}}}
\newcommand{\mehat}[1]{{\hat{m}^{(#1)}}}


% Figure reference, lower-case.
\def\figref#1{figure~\ref{#1}}
% Figure reference, capital. For start of sentence
\def\Figref#1{Figure~\ref{#1}}
\def\twofigref#1#2{figures \ref{#1} and \ref{#2}}
\def\quadfigref#1#2#3#4{figures \ref{#1}, \ref{#2}, \ref{#3} and \ref{#4}}
% Section reference, lower-case.
\def\secref#1{section~\ref{#1}}
% Section reference, capital.
\def\Secref#1{Section~\ref{#1}}
% Reference to two sections.
\def\twosecrefs#1#2{sections \ref{#1} and \ref{#2}}
% Reference to three sections.
\def\secrefs#1#2#3{sections \ref{#1}, \ref{#2} and \ref{#3}}
% Reference to an equation, lower-case.
\def\eqref#1{equation~\ref{#1}}
% Reference to an equation, upper case
\def\Eqref#1{Equation~\ref{#1}}
% A raw reference to an equation---avoid using if possible
\def\plaineqref#1{\ref{#1}}
% Reference to a chapter, lower-case.
\def\chapref#1{chapter~\ref{#1}}
% Reference to an equation, upper case.
\def\Chapref#1{Chapter~\ref{#1}}
% Reference to a range of chapters
\def\rangechapref#1#2{chapters\ref{#1}--\ref{#2}}
% Reference to an algorithm, lower-case.
\def\algref#1{algorithm~\ref{#1}}
% Reference to an algorithm, upper case.
\def\Algref#1{Algorithm~\ref{#1}}
\def\twoalgref#1#2{algorithms \ref{#1} and \ref{#2}}
\def\Twoalgref#1#2{Algorithms \ref{#1} and \ref{#2}}
% Reference to a part, lower case
\def\partref#1{part~\ref{#1}}
% Reference to a part, upper case
\def\Partref#1{Part~\ref{#1}}
\def\twopartref#1#2{parts \ref{#1} and \ref{#2}}

\def\ceil#1{\lceil #1 \rceil}
\def\floor#1{\lfloor #1 \rfloor}
\def\1{\bm{1}}
\newcommand{\train}{\mathcal{D}}
\newcommand{\valid}{\mathcal{D_{\mathrm{valid}}}}
\newcommand{\test}{\mathcal{D_{\mathrm{test}}}}

\def\eps{{\epsilon}}


% Random variables
\def\reta{{\textnormal{$\eta$}}}
\def\ra{{\textnormal{a}}}
\def\rb{{\textnormal{b}}}
\def\rc{{\textnormal{c}}}
\def\rd{{\textnormal{d}}}
\def\re{{\textnormal{e}}}
\def\rf{{\textnormal{f}}}
\def\rg{{\textnormal{g}}}
\def\rh{{\textnormal{h}}}
\def\ri{{\textnormal{i}}}
\def\rj{{\textnormal{j}}}
\def\rk{{\textnormal{k}}}
\def\rl{{\textnormal{l}}}
% rm is already a command, just don't name any random variables m
\def\rn{{\textnormal{n}}}
\def\ro{{\textnormal{o}}}
\def\rp{{\textnormal{p}}}
\def\rq{{\textnormal{q}}}
\def\rr{{\textnormal{r}}}
\def\rs{{\textnormal{s}}}
\def\rt{{\textnormal{t}}}
\def\ru{{\textnormal{u}}}
\def\rv{{\textnormal{v}}}
\def\rw{{\textnormal{w}}}
\def\rx{{\textnormal{x}}}
\def\ry{{\textnormal{y}}}
\def\rz{{\textnormal{z}}}

% Random vectors
\def\rvepsilon{{\mathbf{\epsilon}}}
\def\rvtheta{{\mathbf{\theta}}}
\def\rva{{\mathbf{a}}}
\def\rvb{{\mathbf{b}}}
\def\rvc{{\mathbf{c}}}
\def\rvd{{\mathbf{d}}}
\def\rve{{\mathbf{e}}}
\def\rvf{{\mathbf{f}}}
\def\rvg{{\mathbf{g}}}
\def\rvh{{\mathbf{h}}}
\def\rvu{{\mathbf{i}}}
\def\rvj{{\mathbf{j}}}
\def\rvk{{\mathbf{k}}}
\def\rvl{{\mathbf{l}}}
\def\rvm{{\mathbf{m}}}
\def\rvn{{\mathbf{n}}}
\def\rvo{{\mathbf{o}}}
\def\rvp{{\mathbf{p}}}
\def\rvq{{\mathbf{q}}}
\def\rvr{{\mathbf{r}}}
\def\rvs{{\mathbf{s}}}
\def\rvt{{\mathbf{t}}}
\def\rvu{{\mathbf{u}}}
\def\rvv{{\mathbf{v}}}
\def\rvw{{\mathbf{w}}}
\def\rvx{{\mathbf{x}}}
\def\rvy{{\mathbf{y}}}
\def\rvz{{\mathbf{z}}}

% Elements of random vectors
\def\erva{{\textnormal{a}}}
\def\ervb{{\textnormal{b}}}
\def\ervc{{\textnormal{c}}}
\def\ervd{{\textnormal{d}}}
\def\erve{{\textnormal{e}}}
\def\ervf{{\textnormal{f}}}
\def\ervg{{\textnormal{g}}}
\def\ervh{{\textnormal{h}}}
\def\ervi{{\textnormal{i}}}
\def\ervj{{\textnormal{j}}}
\def\ervk{{\textnormal{k}}}
\def\ervl{{\textnormal{l}}}
\def\ervm{{\textnormal{m}}}
\def\ervn{{\textnormal{n}}}
\def\ervo{{\textnormal{o}}}
\def\ervp{{\textnormal{p}}}
\def\ervq{{\textnormal{q}}}
\def\ervr{{\textnormal{r}}}
\def\ervs{{\textnormal{s}}}
\def\ervt{{\textnormal{t}}}
\def\ervu{{\textnormal{u}}}
\def\ervv{{\textnormal{v}}}
\def\ervw{{\textnormal{w}}}
\def\ervx{{\textnormal{x}}}
\def\ervy{{\textnormal{y}}}
\def\ervz{{\textnormal{z}}}

% Random matrices
\def\rmA{{\mathbf{A}}}
\def\rmB{{\mathbf{B}}}
\def\rmC{{\mathbf{C}}}
\def\rmD{{\mathbf{D}}}
\def\rmE{{\mathbf{E}}}
\def\rmF{{\mathbf{F}}}
\def\rmG{{\mathbf{G}}}
\def\rmH{{\mathbf{H}}}
\def\rmI{{\mathbf{I}}}
\def\rmJ{{\mathbf{J}}}
\def\rmK{{\mathbf{K}}}
\def\rmL{{\mathbf{L}}}
\def\rmM{{\mathbf{M}}}
\def\rmN{{\mathbf{N}}}
\def\rmO{{\mathbf{O}}}
\def\rmP{{\mathbf{P}}}
\def\rmQ{{\mathbf{Q}}}
\def\rmR{{\mathbf{R}}}
\def\rmS{{\mathbf{S}}}
\def\rmT{{\mathbf{T}}}
\def\rmU{{\mathbf{U}}}
\def\rmV{{\mathbf{V}}}
\def\rmW{{\mathbf{W}}}
\def\rmX{{\mathbf{X}}}
\def\rmY{{\mathbf{Y}}}
\def\rmZ{{\mathbf{Z}}}

% Elements of random matrices
\def\ermA{{\textnormal{A}}}
\def\ermB{{\textnormal{B}}}
\def\ermC{{\textnormal{C}}}
\def\ermD{{\textnormal{D}}}
\def\ermE{{\textnormal{E}}}
\def\ermF{{\textnormal{F}}}
\def\ermG{{\textnormal{G}}}
\def\ermH{{\textnormal{H}}}
\def\ermI{{\textnormal{I}}}
\def\ermJ{{\textnormal{J}}}
\def\ermK{{\textnormal{K}}}
\def\ermL{{\textnormal{L}}}
\def\ermM{{\textnormal{M}}}
\def\ermN{{\textnormal{N}}}
\def\ermO{{\textnormal{O}}}
\def\ermP{{\textnormal{P}}}
\def\ermQ{{\textnormal{Q}}}
\def\ermR{{\textnormal{R}}}
\def\ermS{{\textnormal{S}}}
\def\ermT{{\textnormal{T}}}
\def\ermU{{\textnormal{U}}}
\def\ermV{{\textnormal{V}}}
\def\ermW{{\textnormal{W}}}
\def\ermX{{\textnormal{X}}}
\def\ermY{{\textnormal{Y}}}
\def\ermZ{{\textnormal{Z}}}

% Vectors
\def\vzero{{\bm{0}}}
\def\vone{{\bm{1}}}
\def\vmu{{\bm{\mu}}}
\def\vtheta{{\bm{\theta}}}
\def\va{{\bm{a}}}
\def\vb{{\bm{b}}}
\def\vc{{\bm{c}}}
\def\vd{{\bm{d}}}
\def\ve{{\bm{e}}}
\def\vf{{\bm{f}}}
\def\vg{{\bm{g}}}
\def\vh{{\bm{h}}}
\def\vi{{\bm{i}}}
\def\vj{{\bm{j}}}
\def\vk{{\bm{k}}}
\def\vl{{\bm{l}}}
\def\vm{{\bm{m}}}
\def\vn{{\bm{n}}}
\def\vo{{\bm{o}}}
\def\vp{{\bm{p}}}
\def\vq{{\bm{q}}}
\def\vr{{\bm{r}}}
\def\vs{{\bm{s}}}
\def\vt{{\bm{t}}}
\def\vu{{\bm{u}}}
\def\vv{{\bm{v}}}
\def\vw{{\bm{w}}}
\def\vx{{\bm{x}}}
\def\vy{{\bm{y}}}
\def\vz{{\bm{z}}}

% Elements of vectors
\def\evalpha{{\alpha}}
\def\evbeta{{\beta}}
\def\evepsilon{{\epsilon}}
\def\evlambda{{\lambda}}
\def\evomega{{\omega}}
\def\evmu{{\mu}}
\def\evpsi{{\psi}}
\def\evsigma{{\sigma}}
\def\evtheta{{\theta}}
\def\eva{{a}}
\def\evb{{b}}
\def\evc{{c}}
\def\evd{{d}}
\def\eve{{e}}
\def\evf{{f}}
\def\evg{{g}}
\def\evh{{h}}
\def\evi{{i}}
\def\evj{{j}}
\def\evk{{k}}
\def\evl{{l}}
\def\evm{{m}}
\def\evn{{n}}
\def\evo{{o}}
\def\evp{{p}}
\def\evq{{q}}
\def\evr{{r}}
\def\evs{{s}}
\def\evt{{t}}
\def\evu{{u}}
\def\evv{{v}}
\def\evw{{w}}
\def\evx{{x}}
\def\evy{{y}}
\def\evz{{z}}

% Matrix
\def\mA{{\bm{A}}}
\def\mB{{\bm{B}}}
\def\mC{{\bm{C}}}
\def\mD{{\bm{D}}}
\def\mE{{\bm{E}}}
\def\mF{{\bm{F}}}
\def\mG{{\bm{G}}}
\def\mH{{\bm{H}}}
\def\mI{{\bm{I}}}
\def\mJ{{\bm{J}}}
\def\mK{{\bm{K}}}
\def\mL{{\bm{L}}}
\def\mM{{\bm{M}}}
\def\mN{{\bm{N}}}
\def\mO{{\bm{O}}}
\def\mP{{\bm{P}}}
\def\mQ{{\bm{Q}}}
\def\mR{{\bm{R}}}
\def\mS{{\bm{S}}}
\def\mT{{\bm{T}}}
\def\mU{{\bm{U}}}
\def\mV{{\bm{V}}}
\def\mW{{\bm{W}}}
\def\mX{{\bm{X}}}
\def\mY{{\bm{Y}}}
\def\mZ{{\bm{Z}}}
\def\mBeta{{\bm{\beta}}}
\def\mPhi{{\bm{\Phi}}}
\def\mLambda{{\bm{\Lambda}}}
\def\mSigma{{\bm{\Sigma}}}

% Tensor
\DeclareMathAlphabet{\mathsfit}{\encodingdefault}{\sfdefault}{m}{sl}
\SetMathAlphabet{\mathsfit}{bold}{\encodingdefault}{\sfdefault}{bx}{n}
\newcommand{\tens}[1]{\bm{\mathsfit{#1}}}
\def\tA{{\tens{A}}}
\def\tB{{\tens{B}}}
\def\tC{{\tens{C}}}
\def\tD{{\tens{D}}}
\def\tE{{\tens{E}}}
\def\tF{{\tens{F}}}
\def\tG{{\tens{G}}}
\def\tH{{\tens{H}}}
\def\tI{{\tens{I}}}
\def\tJ{{\tens{J}}}
\def\tK{{\tens{K}}}
\def\tL{{\tens{L}}}
\def\tM{{\tens{M}}}
\def\tN{{\tens{N}}}
\def\tO{{\tens{O}}}
\def\tP{{\tens{P}}}
\def\tQ{{\tens{Q}}}
\def\tR{{\tens{R}}}
\def\tS{{\tens{S}}}
\def\tT{{\tens{T}}}
\def\tU{{\tens{U}}}
\def\tV{{\tens{V}}}
\def\tW{{\tens{W}}}
\def\tX{{\tens{X}}}
\def\tY{{\tens{Y}}}
\def\tZ{{\tens{Z}}}


\newcommand{\cA}{{\cal A}}
\newcommand{\cB}{{\cal B}}
\newcommand{\cC}{{\cal C}}
\newcommand{\cD}{{\cal D}}
\newcommand{\cE}{{\cal E}}
\newcommand{\cF}{{\cal F}}
\newcommand{\cG}{{\cal G}}
\newcommand{\cH}{{\cal H}}
\newcommand{\cI}{{\cal I}}
\newcommand{\cJ}{{\cal J}}
\newcommand{\cK}{{\cal K}}
\newcommand{\cL}{{\cal L}}
\newcommand{\cM}{{\cal M}}
\newcommand{\cN}{{\cal N}}
\newcommand{\cO}{{\cal O}}
\newcommand{\cP}{{\cal P}}
\newcommand{\cQ}{{\cal Q}}
\newcommand{\cR}{{\cal R}}
\newcommand{\cS}{{\cal S}}
\newcommand{\cT}{{\cal T}}
\newcommand{\cU}{{\cal U}}
\newcommand{\cV}{{\cal V}}
\newcommand{\cW}{{\cal W}}
\newcommand{\cX}{{\cal X}}
\newcommand{\cY}{{\cal Y}}
\newcommand{\cZ}{{\cal Z}}

% Graph
\def\gA{{\mathcal{A}}}
\def\gB{{\mathcal{B}}}
\def\gC{{\mathcal{C}}}
\def\gD{{\mathcal{D}}}
\def\gE{{\mathcal{E}}}
\def\gF{{\mathcal{F}}}
\def\gG{{\mathcal{G}}}
\def\gH{{\mathcal{H}}}
\def\gI{{\mathcal{I}}}
\def\gJ{{\mathcal{J}}}
\def\gK{{\mathcal{K}}}
\def\gL{{\mathcal{L}}}
\def\gM{{\mathcal{M}}}
\def\gN{{\mathcal{N}}}
\def\gO{{\mathcal{O}}}
\def\gP{{\mathcal{P}}}
\def\gQ{{\mathcal{Q}}}
\def\gR{{\mathcal{R}}}
\def\gS{{\mathcal{S}}}
\def\gT{{\mathcal{T}}}
\def\gU{{\mathcal{U}}}
\def\gV{{\mathcal{V}}}
\def\gW{{\mathcal{W}}}
\def\gX{{\mathcal{X}}}
\def\gY{{\mathcal{Y}}}
\def\gZ{{\mathcal{Z}}}

% Sets
\def\sA{{\mathbb{A}}}
\def\sB{{\mathbb{B}}}
\def\sC{{\mathbb{C}}}
\def\sD{{\mathbb{D}}}
% Don't use a set called E, because this would be the same as our symbol
% for expectation.
\def\sF{{\mathbb{F}}}
\def\sG{{\mathbb{G}}}
\def\sH{{\mathbb{H}}}
\def\sI{{\mathbb{I}}}
\def\sJ{{\mathbb{J}}}
\def\sK{{\mathbb{K}}}
\def\sL{{\mathbb{L}}}
\def\sM{{\mathbb{M}}}
\def\sN{{\mathbb{N}}}
\def\sO{{\mathbb{O}}}
\def\sP{{\mathbb{P}}}
\def\sQ{{\mathbb{Q}}}
\def\sR{{\mathbb{R}}}
\def\sS{{\mathbb{S}}}
\def\sT{{\mathbb{T}}}
\def\sU{{\mathbb{U}}}
\def\sV{{\mathbb{V}}}
\def\sW{{\mathbb{W}}}
\def\sX{{\mathbb{X}}}
\def\sY{{\mathbb{Y}}}
\def\sZ{{\mathbb{Z}}}


% Entries of a matrix
\def\emLambda{{\Lambda}}
\def\emA{{A}}
\def\emB{{B}}
\def\emC{{C}}
\def\emD{{D}}
\def\emE{{E}}
\def\emF{{F}}
\def\emG{{G}}
\def\emH{{H}}
\def\emI{{I}}
\def\emJ{{J}}
\def\emK{{K}}
\def\emL{{L}}
\def\emM{{M}}
\def\emN{{N}}
\def\emO{{O}}
\def\emP{{P}}
\def\emQ{{Q}}
\def\emR{{R}}
\def\emS{{S}}
\def\emT{{T}}
\def\emU{{U}}
\def\emV{{V}}
\def\emW{{W}}
\def\emX{{X}}
\def\emY{{Y}}
\def\emZ{{Z}}
\def\emSigma{{\Sigma}}

% entries of a tensor
% Same font as tensor, without \bm wrapper
\newcommand{\etens}[1]{\mathsfit{#1}}
\def\etLambda{{\etens{\Lambda}}}
\def\etA{{\etens{A}}}
\def\etB{{\etens{B}}}
\def\etC{{\etens{C}}}
\def\etD{{\etens{D}}}
\def\etE{{\etens{E}}}
\def\etF{{\etens{F}}}
\def\etG{{\etens{G}}}
\def\etH{{\etens{H}}}
\def\etI{{\etens{I}}}
\def\etJ{{\etens{J}}}
\def\etK{{\etens{K}}}
\def\etL{{\etens{L}}}
\def\etM{{\etens{M}}}
\def\etN{{\etens{N}}}
\def\etO{{\etens{O}}}
\def\etP{{\etens{P}}}
\def\etQ{{\etens{Q}}}
\def\etR{{\etens{R}}}
\def\etS{{\etens{S}}}
\def\etT{{\etens{T}}}
\def\etU{{\etens{U}}}
\def\etV{{\etens{V}}}
\def\etW{{\etens{W}}}
\def\etX{{\etens{X}}}
\def\etY{{\etens{Y}}}
\def\etZ{{\etens{Z}}}


\newcommand{\Expec}[1]{\mathbb{E}\left[#1\right]}
\newcommand{\suchthat}{\;\ifnum\currentgrouptype=16 \middle\fi|\;}
\newcommand{\fE}{\mathbb{E}}
\renewcommand{\Var}[1]{\mathrm{Var}\left[#1\right]}

\excludeversion{problem}
\includeversion{solution}

\begin{document}

\begin{center}
{\bf Data Privacy: Homeworks \#3 \begin{solution}Solutions\end{solution}}\\ \medskip
\begin{problem}Due on Wednesday 11PM, May 27, 2025 \end{problem}
\end{center}

\begin{problem}
\noindent 
Note: Submit homework via LearnUS.
% Also, we have some extra problems in Problem 7 (3-5).
\LeftFoot{Homework 3}
\end{problem}

\begin{solution}
\LeftFoot{Homework 3 Solution}
\end{solution}

\begin{enumerate}

%%%%%%%%%%%%%%%%%%%%%%%%%%%%% Problem 1 %%%%%%%%%%%%%%%%%%%%%%%%%%%%%
\item  
Show that the minimum mean square estimate is $\hat{X}(Y) = \Expec{X|Y}$.
You can start from
\begin{align*}
    \Expec{(\hat{X}(Y)-X)^2} = \Expec{((\hat{X}(Y)-\Expec{X|Y}) +(\Expec{X|Y}-X) )^2}
\end{align*}

%%%%%%%%%%%%%%%%%%%%%%%% Problem 1 Solution %%%%%%%%%%%%%%%%%%%%%%%%%%
\begin{solution}
{\bf Solution:}\\
    \begin{align*}
        \Expec{(\hat{X}(Y)-X)^2} =& \Expec{((\hat{X}(Y)-\Expec{X|Y}) +(\Expec{X|Y}-X) )^2}\\
         =& \Expec{(\hat{X}(Y)-\Expec{X|Y})^2}   + 2\Expec{(\hat{X}(Y) - \Expec{X|Y})(\Expec{X|Y}-X)} \nonumber\\
         &+ \Expec{(\Expec{X|Y}-X)^2}\\
         =& \Expec{(\hat{X}(Y)-\Expec{X|Y})^2}   + 2\Expec{\Expec{(\hat{X}(Y) - \Expec{X|Y})(\Expec{X|Y}-X)\suchthat Y}}\nonumber\\
         &+ \Expec{(\Expec{X|Y}-X)^2}\\
         =& \Expec{(\hat{X}(Y)-\Expec{X|Y})^2}   + 2\Expec{(\hat{X}(Y) - \Expec{X|Y})(\Expec{X|Y}-\Expec{X|Y}) }\nonumber\\
         &+ \Expec{(\Expec{X|Y}-X)^2}\\
         =& \Expec{(\hat{X}(Y)-\Expec{X|Y})^2}  + \Expec{(\Expec{X|Y}-X)^2}\\
         \geq & \Expec{(\Expec{X|Y}-X)^2} .
    \end{align*}
    The equality holds iff $\hat{X}(Y) = \Expec{X|Y}$ almost surely.
\end{solution}

\vskip .2in

%%%%%%%%%%%%%%%%%%%%%%%%%%%%% Problem 2 %%%%%%%%%%%%%%%%%%%%%%%%%%%%%
\item
Let $X$ be a random variable and $Z~\sim\cN(0, \sigma^2)$ be an independent Gaussian noise.
Let $Y= X+Z$ be noise corrupted version of $X$.
\begin{enumerate}
    \item Show that
\begin{align*}
    \fE[X|Y=y]  = y + \sigma^2 \frac{\partial}{\partial y} \log P_{Y}(y).
\end{align*}
Hint: You can start from 
\begin{align*}
    P_Y(y) = \int P_X(x) P_Z(y-x)\, dx 
\end{align*}
and compare it with $\frac{\partial}{\partial y} P_Y(y)$.

\item {\bf (Bonus)} What about $\Var{X|Y=y}$? 
\end{enumerate}

%%%%%%%%%%%%%%%%%%%%%%%% Problem 2 Solution %%%%%%%%%%%%%%%%%%%%%%%%%%
\begin{solution}
\textbf{Solution:}\\
\begin{enumerate}
    \item
    \begin{align*}
        \partial_y P_Y(y) 
        &= \int P_X(x) \partial_y P_Z(y - x)\, dx \\
        &= \int P_X(x) \partial_y \frac{1}{\sqrt{2\pi\sigma^2}} \exp\left(-\frac{(x - y)^2}{2\sigma^2}\right)\, dx \\
        &= \int P_X(x) \frac{1}{\sqrt{2\pi\sigma^2}} \left(\frac{x - y}{\sigma^2}\right)\exp\left(-\frac{(x - y)^2}{2\sigma^2}\right)\, dx \\
        &= -\frac{y}{\sigma^2} P_Y(y) + \int P_X(x) \frac{1}{\sqrt{2\pi\sigma^2}} \left(\frac{x}{\sigma^2}\right)\exp\left(-\frac{(x - y)^2}{2\sigma^2}\right)\, dx \\
        &= -\frac{y}{\sigma^2} P_Y(y) + \int \left(\frac{x}{\sigma^2}\right) P_X(x) P_{Y|X}(y | x)\, dx \\
        &= -\frac{y}{\sigma^2} P_Y(y) + \int \left(\frac{x}{\sigma^2}\right) P_Y(y) P_{X|Y}(x | y)\, dx \\
        &= -\frac{y}{\sigma^2} P_Y(y) + P_Y(y) \frac{1}{\sigma^2} \fE[X | Y = y]
    \end{align*}
    Since $\partial_y \log P_Y(y) = \frac{\partial_y P_Y(y)}{P_Y(y)}$, we have
    \begin{align*}
        \fE[X | Y = y] = y + \sigma^2 \partial_y \log P_Y(y).
    \end{align*}

    \item
    \begin{align*}
        \partial_y^2 \log P_Y(y)
        &= \frac{P_Y(y)\partial_y^2 P_Y(y)-(\partial_y P_Y(y))^2}{P_Y(y)^2} \\
        &= \frac{\partial_y^2 P_Y(y)}{P_Y(y)} - \left(\frac{\partial_y P_Y(y)}{P_Y(y)}\right)^2 \\
        &= \frac{\partial_y^2 P_Y(y)}{P_Y(y)} - \left(\partial_y \log P_Y(y)\right)^2
    \end{align*}
    Compute the second derivative:
    \begin{align*}
        \partial_y^2 P_Y(y) 
        &= -\partial_y \left(\frac{y}{\sigma^2} P_Y(y)\right) + \partial_y \left(P_Y(y) \cdot \frac{1}{\sigma^2} \fE[X | Y = y]\right) \\
        &= -\frac{1}{\sigma^2} P_Y(y) - \frac{y}{\sigma^2} \partial_y P_Y(y) + \partial_y \int P_X(x) \frac{x}{\sigma^2} P_Z(y - x)\, dx \\
        &= -\frac{1}{\sigma^2} P_Y(y) - \frac{y}{\sigma^2} \partial_y P_Y(y) + \int \frac{x(x - y)}{\sigma^4} P_{X|Y}(x | y) P_Y(y)\, dx \\
        &= -\frac{1}{\sigma^2} P_Y(y) - \frac{y}{\sigma^2} \left(-\frac{y}{\sigma^2} P_Y(y) + P_Y(y) \frac{1}{\sigma^2} \fE[X | Y = y]\right) \\
        &\quad + \frac{1}{\sigma^4} \fE[X^2 | Y = y] P_Y(y) - \frac{y}{\sigma^4} \fE[X | Y = y] P_Y(y)
    \end{align*}
    Dividing both sides by $P_Y(y)$ gives
    \begin{align*}
        \frac{\partial_y^2 P_Y(y)}{P_Y(y)} 
        = \frac{1}{\sigma^4} \left(-\sigma^2 + y^2 - 2y \fE[X | Y = y] + \fE[X^2 | Y = y]\right).
    \end{align*}
    The previous result gives
    \begin{align*}
        \left(\partial_y \log P_Y(y)\right)^2 
        = \frac{1}{\sigma^4} \left( \fE[X | Y = y]^2 + y^2 - 2y \fE[X | Y = y] \right)
    \end{align*}
    Therefore,
    \begin{align*}
        \sigma^4 \partial_y^2 \log P_Y(y) 
        &= \fE[X^2 | Y = y] - \fE[X | Y = y]^2 - \sigma^2 \\
        &= \Var{X | Y = y} - \sigma^2
    \end{align*}
    which implies
    \begin{align*}
        \Var{X | Y = y} = \sigma^2 + \sigma^4 \partial_y^2 \log P_Y(y).
    \end{align*}
\end{enumerate}
\end{solution}


\vskip .2in

%%%%%%%%%%%%%%%%%%%%%%%%%%%%% Problem 3 %%%%%%%%%%%%%%%%%%%%%%%%%%%%%
\item 
Let $q$ be the forwarding distribution of DDPM. 
Show that 
\begin{align*}
    q(X^{(n-1)}|X^{(n)}, X^{(0)}) = \cN\left(\mu_n(X^{(n)}|X^{(0)}), \frac{1-\bar{\alpha}_{n-1}}{1-\bar{\alpha}_n}\beta_n\right)
\end{align*}
where
\begin{align*}
    \mu_n(X^{(n)}|X^{(0)}) = \frac{1}{\sqrt{1-\beta_n}} \left(X^{(n)} + \beta_n \frac{\partial}{\partial X^{(n)}} \log q (X^{(n)}|X^{(0)})\right)
\end{align*}
Hint: Given $X^{(0)}$, all distributions are Gaussian. More precisely,
\begin{align*}
    \left.\begin{bmatrix}
        X^{(n-1)}\\X^{(n)} 
    \end{bmatrix}
    \right|X^{0)}
\end{align*}
 is multivariate Gaussian.
 You can use the fact that if
\begin{align*}
    \begin{bmatrix}
        U\\V
    \end{bmatrix}
    \sim \cN\left(\begin{bmatrix}\mu_U\\\mu_V\\\end{bmatrix},
    \begin{bmatrix}\sigma_{UU}  & \sigma_{UV} \\ \sigma_{VU} & \sigma_{VV}\end{bmatrix}\right)
\end{align*}
then
\begin{align*}
    U|V \sim\cN\left(\mu_U +(\sigma_{UV}/\sigma_{VV}) (V - \mu_V), \sigma_{UU} - \sigma_{UV}^2/\sigma_{VV}\right).
\end{align*}
Also note that we have an analytic formula for
\begin{align*}
    \frac{\partial}{\partial X^{(n)}} \log q (X^{(n)}|X^{(0)})
\end{align*}

%%%%%%%%%%%%%%%%%%%%%%%% Problem 3 Solution %%%%%%%%%%%%%%%%%%%%%%%%%%
\begin{solution}
{\bf Solution:}\\
Since $X^{(n-1)}|X^{(n)}, X^{(0)}$ is Gaussian, it is enough to compute the mean and variance.
Clearly,
    \begin{align*}
        \left.\begin{bmatrix}
            X^{(n-1)}\\X^{(n)} 
        \end{bmatrix}
        \right| X^{(0)}
        \sim \cN\left(\begin{bmatrix}\sqrt{\bar{\alpha}_{n-1}}X^{(0)} \\ \sqrt{\bar{\alpha}_{n}}X^{(0)} \\\end{bmatrix},
        \begin{bmatrix}1-\bar{\alpha}_{n-1} & \sqrt{1-\beta_n} (1-\bar{\alpha}_{n-1}) \\
        \sqrt{1-\beta_n} (1-\bar{\alpha}_{n-1}) & 1-\bar{\alpha}_n\end{bmatrix}\right)
    \end{align*}
Then, the conditional mean is
\begin{align*}
    &\sqrt{\bar{\alpha}_{n-1}} X^{(0)} +\frac{\sqrt{1-\beta_n}(1-\bar{\alpha}_{n-1})}{1-\bar{\alpha}_n} (X^{(n)} - \sqrt{\bar{\alpha}_n}X^{(0)})\nonumber\\
    =& \left(\sqrt{\bar{\alpha}_{n-1}} - \frac{\sqrt{1-\beta_n}(1-\bar{\alpha}_{n-1})}{1-\bar{\alpha}_n} \sqrt{\bar{\alpha}_n}\right) X^{(0)}
    +\frac{\sqrt{1-\beta_n}(1-\bar{\alpha}_{n-1})}{1-\bar{\alpha}_n} X^{(n)} .
\end{align*}
Note that 
\begin{align*}
    & \sqrt{\bar{\alpha}_{n-1}} - \frac{\sqrt{1-\beta_n}(1-\bar{\alpha}_{n-1})}{1-\bar{\alpha}_n} \sqrt{\bar{\alpha}_n}\nonumber\\
    =& \frac{\sqrt{\bar{\alpha}_{n-1}}}{1-\bar{\alpha}_n} (1-\bar{\alpha}_n - (1-\beta_n) (1-\bar{\alpha}_{n-1}))\\
    =& \frac{\sqrt{\bar{\alpha}_{n-1}}}{1-\bar{\alpha}_n} (1-\bar{\alpha}_{n-1}(1-\beta_n) - (1-\beta_n) (1-\bar{\alpha}_{n-1}))\\
    =& \frac{\sqrt{\bar{\alpha}_{n-1}}}{1-\bar{\alpha}_n} \beta_n.
\end{align*}
Thus, the conditional mean is
\begin{align*}
    &\sqrt{\bar{\alpha}_{n-1}} X^{(0)} +\frac{\sqrt{1-\beta_n}(1-\bar{\alpha}_{n-1})}{1-\bar{\alpha}_n} (X^{(n)} - \sqrt{\bar{\alpha}_n}X^{(0)})\nonumber\\
    =& \frac{\sqrt{\bar{\alpha}_{n-1}}}{1-\bar{\alpha}_n} \beta_n X^{(0)}
    +\frac{\sqrt{1-\beta_n}(1-\bar{\alpha}_{n-1})}{1-\bar{\alpha}_n} X^{(n)}.
\end{align*}

Recall $\bar{\alpha}_n = \bar{\alpha}_{n-1}(1-\beta_n)$.
For comparison, 
\begin{align*}
    \mu_n(X^{(n)}|X^{(0)}) =& \frac{1}{\sqrt{1-\beta_n}} \left(X^{(n)} + \beta_n \partial_{X^{(n)}} \log q(X^{(0)}|X^{(0)})\right)\\
    =& \frac{1}{\sqrt{1-\beta_n}} \left(X^{(n)} + \beta_n \left(-\frac{X^{(n)}-\sqrt{\bar{\alpha}_n}X^{(0)}}{1-\bar{\alpha}_n}\right)\right)\\
    =&\frac{\sqrt{\bar{\alpha}_{n-1}}}{1-\bar{\alpha}_n} \beta_n X^{(0)} 
    +\frac{1}{\sqrt{1-\beta_n}(1-\bar{\alpha}_n)} \left( 1-\bar{\alpha}_n- \beta_n\right)X^{(n)}\\
    =&\frac{\sqrt{\bar{\alpha}_{n-1}}}{1-\bar{\alpha}_n} \beta_n X^{(0)} 
    +\frac{1-\bar{\alpha}_{n-1}}{(1-\bar{\alpha}_n)} \sqrt{1-\beta_n}X^{(n)}.
\end{align*}


On the other hand, the conditional variance is
\begin{align*}
    & 1-\bar{\alpha}_{n-1} - \frac{(1-\beta_n)(1-\bar{\alpha}_{n-1})^2}{1-\bar{\alpha}_n}\nonumber\\
    &= \frac{1-\bar{\alpha}_{n-1}}{1-\bar{\alpha}_n} \left(1-\bar{\alpha}_n- (1-\beta_n)(1-\bar{\alpha}_{n-1})\right)\\
    &= \frac{1-\bar{\alpha}_{n-1}}{1-\bar{\alpha}_n} \left(1-\bar{\alpha}_{n-1}(1-\beta_n) - (1-\beta_n)(1-\bar{\alpha}_{n-1})\right)\\
    &= \frac{1-\bar{\alpha}_{n-1}}{1-\bar{\alpha}_n} \beta_n
\end{align*}
\end{solution}

\vskip .2in
%%%%%%%%%%%%%%%%%%%%%%%%%%%%% Problem 4 %%%%%%%%%%%%%%%%%%%%%%%%%%%%%
\item 
Show that
\begin{align*}
    D_{\mathrm{KL}}\left(\cN(\mu_0, \sigma_0^2I) \|\cN(\mu_1, \sigma_1^2I) \right)
        = \frac{1}{2\sigma_1^2} \|\mu_1-\mu_0\|^2 + \frac{d}{2}\left(\frac{\sigma_0^2}{\sigma_1^2}-1\right) + d\log \frac{\sigma_1}{\sigma_0}.
\end{align*}
where $d$ is dimension of Gaussian distributions and $\mu_0, \mu_1\in\fR^d, \sigma_1, \sigma_2>0$.

%%%%%%%%%%%%%%%%%%%%%%%% Problem 4 Solution %%%%%%%%%%%%%%%%%%%%%%%%%%
\begin{solution}
{\bf Solution:}\\
\begin{align*}
        &D_{KL}(\cN(\mu_0, \sigma_0^2I) \|\cN(\mu_1, \sigma_1^2I) )\nonumber\\
        &=\fE_{\mu_0, \sigma_0}\left[\log\frac{\frac{1}{\sqrt{(2\pi \sigma_0)^d}} \exp\left(-\frac{1}{2\sigma_0^2}\|X-\mu_0\|^2\right)}{\frac{1}{\sqrt{(2\pi \sigma_1)^d}} \exp\left(-\frac{1}{2\sigma_1^2}\|X-\mu_1\|^2\right)}\right]\\
        &=d\log \frac{\sigma_1}{\sigma_0} + \fE_{\mu_0, \sigma_0}\left[-\frac{1}{2\sigma_0^2}\|X-\mu_0\|^2+\frac{1}{2\sigma_1^2}\|X-\mu_1\|^2\right]
\end{align*}
Since $\fE_{\mu_0, \sigma_0}[X-\mu_0] = 0$, 
the second term is simply
\begin{align*}
        &\fE_{\mu_0, \sigma_0}\left[-\frac{1}{2\sigma_0^2}\|X-\mu_0\|^2+\frac{1}{2\sigma_1^2}\|X-\mu_1\|^2\right]\nonumber\\
        =&\fE_{\mu_0, \sigma_0}\left[-\frac{1}{2\sigma_0^2}\|X-\mu_0\|^2+\frac{1}{2\sigma_1^2}\|X-\mu_0\|^2+\frac{1}{2\sigma_1^2}\|\mu_1-\mu_0\|^2\right]\\
        =&\fE_{\mu_0, \sigma_0}\left[\left(\frac{1}{2\sigma_1^2}-\frac{1}{2\sigma_0^2}\right)\|X-\mu_0\|^2\right]+\frac{1}{2\sigma_1^2}\|\mu_1-\mu_0\|^2\\
        =&\left(\frac{1}{2\sigma_1^2}-\frac{1}{2\sigma_0^2}\right)d\sigma_0^2+\frac{1}{2\sigma_1^2}\|\mu_1-\mu_0\|^2\\
        =&\frac{1}{2\sigma_1^2}\|\mu_1-\mu_0\|^2+\frac{d}{2}\left(\frac{\sigma_0^2}{\sigma_1^2}-1\right).
\end{align*}
\end{solution}

\vskip .2in
%%%%%%%%%%%%%%%%%%%%%%%%%%%%% Problem 6 %%%%%%%%%%%%%%%%%%%%%%%%%%%%%
\item 
Consider the DDIM forward process.
\begin{align}
    q(X^{(N)}|X^{(0)}) =&  \cN(\sqrt{\bar{\alpha}_N} X^{(0)}, (1-\bar{\alpha}_N)I)\nonumber\\
    q(X^{(n)}|X^{(n+1)}, X^{(0)}) =& \cN(\sqrt{\bar{\alpha}_n} X^{(0)} + \frac{\sqrt{1-\bar{\alpha}_n 
    - \sigma_{n+1}^2}}{\sqrt{1-\bar{\alpha}_{n+1}}} (X^{(n+1)}-\sqrt{\bar{\alpha}_{n+1}} X^{(0)}), \sigma_{n+1}^2 I)\label{eq:ddim forward}
\end{align}
for $\sigma_n\geq 0$.
    Show that \eqref{eq:ddim forward} matches the marginal distribution of DDPM.

%%%%%%%%%%%%%%%%%%%%%%%% Problem 6 Solution %%%%%%%%%%%%%%%%%%%%%%%%%%
\begin{solution}
{\bf Solution:}\\
For $n=N$, it is clear that the distribution of $q(X^{(N)}|X^{(0)})$ matches the distribution.
\begin{align*}
    X^{(N)} \stackrel{(d)}{=} \sqrt{\bar{\alpha}_N} X^{(0)} + Z^{(N)}
\end{align*}
where $Z^{(N)}\sim\cN(0, (1-\bar{\alpha}_N^2))$.
Suppose DDIM forward process matches with DDPM at $n+1$, i.e.,
\begin{align*}
    X^{(n+1)}\stackrel{(d)}{=} \sqrt{\bar{\alpha}_{n+1}} X^{(0)} + Z^{(n+1)} 
\end{align*}
where $Z^{(n+1)}\sim\cN(0, (1-\bar{\alpha}_{n+1}^2))$.
Then, we will show that DDIM forward process also matches with DDPM at $n$.
From \eqref{eq:ddim forward},
\begin{align*}
    X^{(n)} \stackrel{(d)}{=} \sqrt{\bar{\alpha}_n} X^{(0)} + \frac{\sqrt{1-\bar{\alpha}_n 
    - \sigma_{n+1}^2}}{\sqrt{1-\bar{\alpha}_{n+1}}} (X^{(n+1)}-\sqrt{\bar{\alpha}_{n+1}} X^{(0)})+ W^{(n+1)}
\end{align*}
where $W^{(n+1)}\sim\cN(0, \sigma_{n+1}^2 I)$.
Then, from assumption where DDPM and DDIM matches at $n+1$, we have
\begin{align*}
    X^{(n)} \stackrel{(d)}{=} \sqrt{\bar{\alpha}_n} X^{(0)} 
    + \frac{\sqrt{1-\bar{\alpha}_n - \sigma_{n+1}^2}}{\sqrt{1-\bar{\alpha}_{n+1}}} Z^{(n+1)} + W^{(n+1)}.
\end{align*}
The sum of two independent noise term has zero mean and covariance of
\begin{align*}
    &\frac{1-\bar{\alpha}_n - \sigma_{n+1}^2}{1-\bar{\alpha}_{n+1}} (1-\bar{\alpha}_{n+1}) I + \sigma_{n+1}^2 I\\
    &= (1-\bar{\alpha}_n)I.
\end{align*}
Thus,
\begin{align*}
    X^{(n)} \sim \cN(\sqrt{\bar{\alpha}_n}X^{(0)}, (1-\bar{\alpha}_n)I).
\end{align*}

\end{solution}

\vskip .2in
%%%%%%%%%%%%%%%%%%%%%%%%%%%%% Problem 8 %%%%%%%%%%%%%%%%%%%%%%%%%%%%%
\item \textbf{Guidance.}  
In diffusion models, classifier-guided conditional generation uses 
\[
\nabla_x \log p_t(x)+\omega\nabla_x \log  p_t(y\,|\,x)
\]
with a classifier scale parameter $\omega\ge 1$.
However, many papers in the literature add classifier guidance to an already conditional model via
\[
\nabla_x \log p_t(x\,|\,y)+s \nabla_x \log p_t(y\,|\,x)
\]
with $s\ge 0$.
Show that the two are equivalent with $s=\omega-1$.

%%%%%%%%%%%%%%%%%%%%%%%% Problem 8 Solution %%%%%%%%%%%%%%%%%%%%%%%%%%
\begin{solution}
{\bf Solution:}\\
By Bayes's theorem, we have $p_t(x\,|\,y)p_t(y)=p_t(y\,|\,x)p_t(x)$. So we have
\begin{align*}
\nabla_x \log p_t(x\,|\,y)+s \nabla_x \log p_t(y\,|\,x)
&=\nabla_x\log\{ p_t(x\,|\,y)p_t(y)\}+s\nabla_x \log p_t(y\,|\,x)\\
&=\nabla_x\log\{ p_t(y\,|\,x)p_t(x)\}+s\nabla_x \log p_t(y\,|\,x)\\
&=\nabla_x \log p_t(x)+(s+1)\nabla_x \log  p_t(y\,|\,x).
\end{align*}
Therefore, with $s=\omega-1$, we can see that the two formula are equivalent.
\end{solution}

\vskip .2in
%%%%%%%%%%%%%%%%%%%%%%%%%%%%% Problem 9 %%%%%%%%%%%%%%%%%%%%%%%%%%%%%
\item \textbf{Langevin}  
Let $p(x)$ is a probability density function that is smooth and strictly positive for all $x\in \mathbb{R}^d$.
Let $\{p_t\}_{t\in[0,T]}$ be the marginal density functions of the Langevin SDE
\[
dX_t=\frac{1}{2}\nabla_{X_t} \log p(X_t)dt+dW_t.
\]
Show that if $p_0=p$, then $p_t=p$ for all $t>0$.

%%%%%%%%%%%%%%%%%%%%%%%% Problem 9 Solution %%%%%%%%%%%%%%%%%%%%%%%%%%
\begin{solution}
{\bf Solution:}\\
The Fokker-Planck equation of Langevin SDE is
\[
\partial_t p_t = -\frac{1}{2} \nabla_{X_t} \cdot (\nabla_{X_t} \log p(X_t) \cdot p_t ) + \frac{1}{2}\Delta p_t.
\]
If $p_t = p$, then
\begin{align*}
    \partial_t p_t &= -\frac{1}{2} \nabla_{X_t} \cdot (\nabla_{X_t}\log p(X_t) \cdot p(X_t)) + \frac{1}{2} \Delta p(X_t)\\
    &= \frac{1}{2} \left( -\nabla_{X_t} \cdot \nabla_{X_t}p(X_t) + \Delta p(X_t)  \right)\\
    &=0
\end{align*}
We have $p_0 = p$ and we know that there is a unique solution. Therefore, the unique solution is $p_t = p$. 
\end{solution}

% \vskip .2in
% %%%%%%%%%%%%%%%%%%%%%%%%%%%%% Problem 11 %%%%%%%%%%%%%%%%%%%%%%%%%%%%%
% \item \textbf{Vanishing} \(t\).
% Consider the 1-dimensional Ornstein--Uhlenbeck process
% \[
% dX_t=-\frac{1}{2}X_t dt+ dW_t
% \]
% for $t\in[0,T]$, where $X_0\sim p_0$.
% % (So $X_t\in \mathbb{R}$.)
% For simplicity, let $p_0=\mathcal{N}\left(0,1\right)$.
% Let
% \[
% \gamma_t=e^{- t/2},
% \qquad
% \sigma_t^2=1-e^{-t}.
% \]
% Consider the loss
% \begin{align*}
% \mathcal{L}(\theta)
% &=
% \mathop{\mathbb{E}}_{t\sim\mathrm{Uniform}([\delta,T])}
% \left[
% \mathop{\mathbb{E}}_{X_0\sim p_0}
% \left[
% \mathop{\mathbb{E}}_{X_t\,|\,X_0}
% \left[ \lambda(t)\left( 
% s_\theta(X_t,t)-\frac{d}{dX_t}\log p_{t|0}(X_t\,|\,X_0)\right)^2 
% \,\bigg|\,X_0\right]
% \right]
% \right]
% \\
% &=
% \mathop{\mathbb{E}}_{\substack{t\sim\mathrm{Uniform}([\delta,T])\\
% X_0\sim p_0\\\varepsilon\sim \mathcal{N}(0,I)}}
% \left[ \frac{\lambda(t)}{\sigma_t^2}\left( \varepsilon_{\theta} (\gamma_t X_0-\sigma_t\varepsilon,t)-\varepsilon\right)^2 \right],
% \end{align*}
% where $\delta\ge 0$, $\lambda(t)\ge 0$ is a continuous function, $s_\theta$ is a score network, and $\varepsilon_\theta(X_t,t)=\sigma_ts_\theta(X_t,t)$.
% It is customary to use $\delta>0$ to ``avoid numerical instabilities.'' 
% In this problem, we explore issues that arise when $\delta=0$.
% \begin{itemize}
% \item[(a)] Show that $p_t=\mathcal{N}\left(0,1\right)$ for all $t>0$.
% \item[(b)]
% Assume $s_\theta$ has been perfectly trained, i.e., $s_\theta(X_t,t)=\frac{d}{dX_t}\log p_t(X_t)=-X_t$.
% Show that if $\min_{t\in[0,T]}\lambda(t)>0$, then
% \begin{align*}
% \mathcal{L}(\theta)&\ge 
% \left(\min_{t\in[0,T]}\lambda(t)\right)
% \mathop{\mathbb{E}}_{t,X_0,X_t}\left[
% \left(s_\theta(X_t,t)-\frac{d}{dX_t}\log p_{t|0}(X_t\,|\,X_0)\right)^2 \right]\\\
% &=\infty.
% \end{align*}
% \item[(c)]
% Show that if $\lambda(t)=\sigma_t^2$ 
% (so $\min_{t\in[0,T]}\lambda(t)=0$)
% and if
% $s_\theta(X_t,t)=\frac{d}{dX_t}\log p_t(X_t)$, then
% \[
% \mathcal{L}(\theta)<\infty.
% \]
% \item[(d)]
% Let $\lambda(t)=\sigma_t^2$.
% Assume there is a $\theta^\star$ such that 
% $s_{\theta^\star}(X_t,t)=\frac{d}{dX_t}\log p_t(X_t)$.
% Let $\theta$ be such that $s_\theta(X_t,t)=\frac{m}{\sigma_t} X_t+\frac{d}{dX_t}\log p_t(X_t)$ for some small $m>0$.
% Show that
% \[
% \mathcal{L}(\theta)-\mathcal{L}(\theta^\star)= m^2.
% \]
% (Conceptually, $m^2$ is small, so $s_\theta$ is nearly optimal with respect to the loss $\mathcal{L}$.)
% \item[(e)]
% Let $s_\theta(X_t,t)=\frac{m}{\sigma_t} X_t+\frac{d}{dX_t}\log p_t(X_t)$ for some small $m>0$.
% Show that the reverse sampling ODE with the trained score $s_\theta$ is of the form 
% \[
% d\overline{X}_t=F(\overline{X}_t,t)dt,
% \]
% where $F(X_t,t)$ blows up as $t\rightarrow 0$.
% (Since the ODE is singular, we expect numerical solutions of it via discretizations to be numerically unstable.)
% \end{itemize}
% 
% \emph{Remark.} The ODE 
% \[
% d\overline{X}_t=-\frac{1}{\sqrt{t}}\overline{X}_t
% \]
% has a general solution $\overline{X}_t=\exp(-2\sqrt{t})$ for $t\ge 0$, so a singular ODE (an ODE with a RHS that blows up) does not necessarily have a singular solution (a solution that blows up).
% 
% %%%%%%%%%%%%%%%%%%%%%%%% Problem 11 Solution %%%%%%%%%%%%%%%%%%%%%%%%%%
% \begin{solution}
% {\bf Solution:}\\
%     \textbf{(a)}
%     By the Fokker--Planck equation, the $\{ p_t \}_{t=0}^T$ satisfies
%     \[
%     \partial_t p_t = - \partial_x( -\frac{x}{2} p_t (x)) + \frac{1}{2} \partial^2_x p_t(x)
%     \]
%     We show that $p_t \sim \mathcal{N}(0,1)$ by verifying that $\mathcal{N}(0,1)$ satisfies above equation
%     \begin{align*}
%     \text{LHS} &= \partial_t p_t = 0 \\
%     \text{RHS} &= -\partial_x \left( -\frac{x}{2\sqrt{2\pi}} \exp{\left( -\frac{1}{2}x^2\right)} \right) + \frac{1}{2} \partial^2_x \left( \frac{1}{\sqrt{2\pi}} \exp{\left( -\frac{1}{2}x^2\right)} \right) \\
%     &= \left(-\frac{x^2}{2\sqrt{2\pi}}  + \frac{1}{2\sqrt{2\pi}} -\frac{1}{2\sqrt{2\pi}}  + \frac{x^2}{2\sqrt{2\pi}} \right)\exp{\left( -\frac{1}{2}x^2\right)} \\
%     &= 0
%     \end{align*}
%     For problem (b)-(e), note that \begin{align*}
%         \mathbb{E}_{X_t,X_0}[X_t X_0] &= \mathbb{E}_{X_0} \left[ \mathbb{E}_{X_t |X_0}\left[ X_t |X_0 \right] X_0 \right]\\
%         &= \mathbb{E}_{X_0}\left[ e^{-t/2}X_0^2\right]\\
%         &= e^{-t/2}
%     \end{align*}
%     
%     
%     \textbf{(b)}
%     Since $p_{t|0} (X_t | X_0) \sim \mathcal{N}\left( e^{-\frac{t}{2}}X_0 , 1-e^{-t} \right)$, we can rewrite $\mathcal{L}(\theta)$ as
%     \begin{align*}
%     \mathcal{L}(\theta) &\ge \left( \min_{t\in[0,T]} \lambda(t) \right) \mathbb{E}_{t,X_0, X_t} \left[ 
%     \left( -X_t + \frac{X_t - e^{-t/2}X_0}{1-e^{-t}} \right)^2
%     \right]\\
%     &= \left( \min_{t\in[0,T]} \lambda(t) \right) \mathbb{E}_{t,X_0, X_t} \left[ \left( \frac{e^{-t}X_t - e^{-t/2}X_0}{1-e^{-t}}  \right)^2 \right] \\
%     & = \left( \min_{t\in[0,T]} \lambda(t) \right)
%     \mathbb{E}_{t \in [\delta, T]} \mathbb{E}_{X_0, X_t \sim \mathcal{N}(0,I)}\left[ A(t) X_t ^2 + B(t) X_t X_0 + C(t) X_0^2 \right]\\
%     &= \left( \min_{t\in[0,T]} \lambda(t) \right)
%     \mathbb{E}_{t \in [\delta, T]} \left[ \frac{e^{-t}-e^{-2t}}{(1-e^{-t})^2}\right]\\
%     &= \left( \min_{t\in[0,T]} \lambda(t) \right)\frac{1}{T-\delta} \int_{\delta}^{T} \left( \frac{1}{e^{t/2} - e^{-t/2}} \right)^2 - \frac{1}{(1+e^{t})^2} dt\\
%     &= \left( \min_{t\in[0,T]} \lambda(t) \right) \frac{1}{T-\delta}\left[ \frac{1}{2-2e^{t}}
%     -t - \frac{1}{e^t +1} + \log (e^t +1)
%     \right]_{\delta}^{T}\\
%     &\rightarrow \infty \quad\text{as}\quad \delta \rightarrow 0
%     \end{align*}
% 
%     \textbf{(c)}
%     \begin{align*}
%     \mathcal{L}(\theta) 
% &= \mathbb{E}_{t,X_0,X_t} \left[ (1-e^{-t})^2 \left( \frac{e^{-t}X_t - e^{-t/2}X_0}{1-e^{-t}}\right)^2 \right]\\
% &= \mathbb{E}_{t,X_0,X_t} \left[ (e^{-t}X_t - e^{-t/2}X_0)^2 \right] \\
% &= \mathbb{E}_{t\in[0,T]} \left[ -e^{-2t} + e^{-t} \right] \\
% &= \frac{1}{T}\left(\frac{1}{2} +\frac{1}{2}e^{-2T}-e^{-T} \right) < \infty
% \end{align*}
% \textbf{(d)}
% \textbf{(sol1)}
% \begin{align*}
% \mathcal{L}(\theta ^\star) &= \mathbb{E}_{t\in[0,T]} \left[ -e^{-2t} + e^{-t} \right]\\
% \mathcal{L}(\theta) &= \mathbb{E}_{t,X_0,X_t} \left[ \sigma_t^2 \left( \frac{m}{\sigma_t} X_t -X_t + \frac{X_t -e^{-t/2}X_0}{\sigma_t} \right)^2 \right]\\
% &= \mathbb{E}_{t,X_0 ,X_t} \left[ (m+e^{-t})X_t -e^{-t/2}X_0 \right]^2 \\
% &= \mathbb{E}_{t} \left[ (m+e^{-t})^2 -2(m+e^{-t})e^{-t}+ e^{-t} \right] \\
% &= m^2 + \mathbb{E}_{t\in[0,T]} \left[ -e^{-2t} + e^{-t} \right]
% \end{align*}
% Then we have $ \mathcal{L}(\theta)- \mathcal{L}(\theta^\star) = m^2 $
% 
% \textbf{(sol2)}
% 
% \begin{align*}
% \mathcal{L}(\theta) &=
% \mathbb{E}
% \left[
% \lambda(t)
% \|\frac{m}{\sigma_t}X_t+\frac{d}{dX_t}\log p_t(X_t)
% -\frac{d}{dX_t}\log p_{t|0}(X_t\,|\,X_0)
% \|^2
% \right]\\
% &=
% \mathbb{E}
% \left[
% \lambda(t)
% \|\frac{d}{dX_t}\log p_t(X_t)
% -\frac{d}{dX_t}\log p_{t|0}(X_t\,|\,X_0)
% \|^2
% \right]
% +
% \mathbb{E}
% \left[
% \lambda(t)
% \|\frac{m}{\sigma_t}X_t
% \|^2
% \right]\\
% & \qquad + \mathbb{E}\left[ 2\lambda(t) \frac{m}{\sigma_t}X_t  ( \frac{d}{dX_t}\log p_t(X_t) -\frac{d}{dX_t}\log p_{t|0}(X_t\,|\,X_0)) \right]\\
% &= \mathcal{L}(\theta^\star) + m^2 + 0.
% \end{align*}
% The last equation is from $\frac{d}{dX_t} \log p_t (X_t) = \mathbb{E}_{X_0 | X_t} \left[ \frac{d}{dX_t} \log p_{t|0} (X_t |X_0 ) | X_t \right]$.
% 
% 
% \textbf{(e)}
% The reverse sampling ODE using $s_\theta$ is 
% \begin{align*}
%     d\bar{X}_t &= \left( -\frac{\bar{X}_t}{2} - \frac{1}{2} (\frac{m}{\sigma_t}X_t -X_t)  \right)dt \\
%     &= -\frac{m}{2\sigma_t}X_t dt
% \end{align*}
% Since $\sigma_t$ goes to zero as $t \rightarrow 0$, the RHS blows up.
% \end{solution}



\end{enumerate}
\end{document}